% Final project proposal, ANLP.
% Build:  cd proposal && pdflatex proposal.tex && pdflatex proposal.tex
% Then rename the output to <TeamNameWithNoSpaces>-Proposal.pdf
\documentclass[10pt,twocolumn,a4paper]{article}

\usepackage[T1]{fontenc}
\usepackage[utf8]{inputenc}
\usepackage{times}
\usepackage[margin=2.0cm,columnsep=0.65cm]{geometry}
\usepackage{booktabs}
\usepackage{microtype}
\usepackage{url}
\usepackage{amsmath}
\usepackage[round,authoryear]{natbib}
\usepackage[hidelinks]{hyperref}

\setlength{\parskip}{0pt}
\setlength{\parindent}{1em}
\frenchspacing

\title{\vspace{-1.2cm}\textbf{Event-Aware Representation Learning for\\
Neutral Multi-Source News Summarization}\vspace{-0.2cm}}

\author{
Chaitanya Thogata \quad Harshit Lalwani \quad Manas Agrawal \quad Raunak Seksaria \\
\small 2023113018 \quad\quad 2023111028 \quad\quad\, 2023113023 \quad\quad\,\, 2023113019 \\
\small IIIT Hyderabad \\
\small \texttt{\{chaitanya.thogata, harshit.lalwani, manas.agrawal, raunak.seksaria\}@research.iiit.ac.in}
}
\date{}

\begin{document}
\maketitle

\section{Project Idea}

We propose a neutral multi-source news summarization system whose internal
representation is explicitly organised around the \emph{event} being reported
rather than the \emph{publisher} reporting it. Given several articles about the
same story from outlets of differing political leaning, the model produces a
single neutral summary; alongside the usual generation objective, its encoder is
trained with a triplet loss that pulls together differently-framed reports of the
same event and pushes apart same-publisher reports of different events.

The deliverable is a summarization model with an event-aware
representation-learning objective, evaluated on how much event content it
retains against how much publisher-specific information remains recoverable from
its representations.

Before writing this proposal we profiled every corpus we intend to use,
reimplemented the evaluation metric the relevant literature reports, and ran the
core experiment once end to end. It works: the objective cuts publisher
recoverability by $63\%$ at unchanged ROUGE, but leaves the framing-bias metric
essentially untouched. That gap is what the project is now about. The findings
below are the substance of our motivation, not background detail.

\section{Motivation}

\paragraph{Nothing constrains what the encoder organises around.}
NeuS \cite{lee2022neus} defines the neutral multi-news summarization task and
shows that generation models can hallucinate politically biased content, not
merely factually incorrect content. Its training signal, however, supervises only
the output text. Independently, \citet{baly2020detect} show that news encoders
readily latch onto source-specific characteristics --- a publisher's stylistic
fingerprint --- rather than the content a task actually needs. If a
summarization encoder does the same, two failure modes follow: \emph{framing
leakage}, where a summary's tone shifts with which outlets happened to be
sampled even though the text reads as neutral, and \emph{neutrality by
deletion}, where the model appears neutral only by dropping contested details
instead of reconciling them.

\paragraph{Publisher identity is present and stance does not transfer.}
We verified the premise rather than assuming it. On the 9{,}282 articles of the
NeuS corpus, a plain TF-IDF logistic probe recovers the publishing outlet at
$17.8\%$ accuracy across $229$ outlets, against a $9.2\%$ majority baseline.
That is roughly twice the corresponding baseline, so outlet identity is plainly
written into the words. Stance behaves in the opposite way. Predicting an
article's left/centre/right leaning reaches only $40.2\%$ when outlets are shared
between train and test, against a $37.3\%$ majority baseline; and when no outlet
appears on both sides it falls to $35.4\%$, \emph{below} the $46.0\%$ obtainable
by always predicting the most frequent class and barely above the $33.3\%$ chance
level. Lexical features therefore encode \emph{which outlet wrote this} and not a
transferable notion of ideology --- the signal a model trained on this data will
reach for first is exactly the one we want its representation to discard. This
reproduces \citeauthor{baly2020detect}'s argument on the very corpus we intend to
generate from, and it is precisely the confound an event-aware objective should
remove.

\paragraph{The field's neutrality metric rewards saying less.}
The framing-bias metric introduced with NeuS scores a generation by summing
NRC-VAD \cite{mohammad2018vad} arousal over the tokens it uses that the neutral
reference does not. We reimplemented it --- the released notebook does not run
--- and measured it on gold data. Across all 9{,}282 gold articles the metric
correlates $0.88$ with the number of calibrated tokens, and once divided by that
count the affect rate is flat: $0.1767$, $0.1763$ and $0.1792$ for left, centre
and right respectively, against a standard deviation of $0.081$. The same
$\approx\!0.177$ appears on WCEP-10 \cite{ghalandari2020wcep}, a different corpus
from a different period. The metric separates partisan from centrist text only
because partisan text diverges more in vocabulary from the reference, not because
that vocabulary carries more affect.

The consequence is direct: a model can lower its framing-bias score simply by
generating fewer novel tokens. The headline metric is blind to neutrality by
deletion --- one of the two failure modes the task exists to prevent. To our
knowledge this has not been reported.

\paragraph{The standard split cannot test generalisation.}
NeuS's released train/validation/test split is random over stories. Recovering
the publisher labels, we find that $92$ of the $99$ test outlets also occur in
training. Any claim about publisher leakage measured on this split is a claim
about outlets the model has already seen.

\section{Preliminary Results}

The data pipeline, the metric, and a first end-to-end experiment are already
built and run. Table~\ref{tab:corpora} summarises what each corpus can support;
Table~\ref{tab:probe} is the result that motivates the rest of the project.

\begin{table}[t]
\centering\small
\setlength{\tabcolsep}{3.5pt}
\begin{tabular}{@{}lccc@{}}
\toprule
\textbf{Corpus} & \textbf{Event} & \textbf{Publisher} & \textbf{Neutral} \\
\midrule
NeuS (3{,}095 stories)     & yes & yes & yes \\
AllSides (37{,}554 arts.)  & no  & yes & no  \\
WCEP-10 (10{,}200 clus.)   & yes & no  & partial \\
WNC (181k sent. pairs)     & no  & no  & partial \\
\bottomrule
\end{tabular}
\caption{An event-aware triplet objective needs all three labels on one example.
Only NeuS supplies them together, and it is the smallest of the four.}
\label{tab:corpora}
\end{table}

\begin{table}[t]
\centering\small
\setlength{\tabcolsep}{3pt}
\begin{tabular}{@{}lrrr@{}}
\toprule
\textbf{Representation} & \textbf{Pub.} & \textbf{Stance} & \textbf{R-1} \\
\midrule
TF-IDF (lexical floor)      & 17.8 & 40.2 & --- \\
BART, generation loss only  & 21.5 & 48.3 & 38.63 \\
BART + event-aware triplet  & \textbf{7.9} & 38.6 & 38.72 \\
\midrule
majority baseline           & 9.2  & 37.3 & --- \\
\bottomrule
\end{tabular}
\caption{Logistic probes on frozen encoder representations of 9{,}282 articles;
publisher is a 229-way problem. Training for summarization makes the publisher
\emph{more} recoverable than raw word counts. The triplet objective drives it
below the majority baseline while ROUGE-1 is unchanged.}
\label{tab:probe}
\end{table}

\paragraph{Setup.} BART-base fine-tuned on the 2{,}452 NeuS training examples for
five epochs at fp16 on a single RTX 2080 Ti, once with the generation loss alone
and once with $\lambda=1.0$, $m=1.0$. Six minutes per run. Probes are logistic
regressions on the pooled encoder state, evaluated on a held-out article split.

\paragraph{Publisher identity is learned, then removed.} A TF-IDF probe recovers
the outlet at $17.8\%$ over 229 classes against a $9.2\%$ majority baseline, so
outlet identity is written into the words. Fine-tuning for summarization makes
this \emph{worse}, not better: the baseline encoder reaches $21.5\%$. The
summarization objective actively organises the representation around the
publisher. Adding the triplet term drops it to $7.9\%$ --- below the majority
baseline, i.e.\ effectively unrecoverable --- a $63\%$ relative reduction, with
ROUGE-1 unchanged at $38.7$ versus $38.6$ and identical mean output length.

This is not representation collapse: a collapsed encoder would put the triplet
loss at the margin, $1.0$, and it converges to $0.047$.

\paragraph{But neutrality does not follow.} The framing-bias metric barely moves
--- $1.732 \rightarrow 1.701$ raw, $0.1774 \rightarrow 0.1751$ per calibrated
token, on identical generation lengths, so this is a genuine non-effect rather
than the length artifact described below. \textbf{Removing publisher information
from the encoder does not by itself make the output more neutral.}
Representation-level leakage and surface-level framing bias are separable, and
the project's central question becomes what closes that gap.

\paragraph{A caveat we will report.} Stance recoverability falls alongside
publisher ($48.3 \rightarrow 38.6$). Because stance is an outlet-level property
in this corpus, the objective cannot strip publisher identity while preserving
ideology --- the inverse of what \citet{baly2020detect} achieve, and a real
limitation of applying their technique in this direction.

\section{Background and Related Work}

\textbf{NeuS} \cite{lee2022neus} introduces the task, the AllSides-derived
corpus of left/centre/right article triples with human-written neutral
``roundup'' summaries, and the framing-bias metric. It fine-tunes BART
\cite{lewis2020bart} on the concatenated sources and reports that a multi-task
auxiliary objective over the encoder helps --- but the auxiliary task is
classification, and nothing ties the representation to event identity.

\textbf{We Can Detect Your Bias} \cite{baly2020detect} predicts article-level
political ideology and uses a triplet objective to strip publisher identity from
a representation while preserving ideology. It also contributes the media-based
split protocol, in which whole outlets are held out. We invert the objective:
preserve \emph{event} identity while removing \emph{source} information.

\textbf{Multi-document event relation graphs} \cite{lei2024multidoc} target the
same framing-bias mitigation goal on the same corpus by building an explicit
event relation graph and reasoning over it with an LLM. Our approach is the
representation-level counterpart: rather than externalising event structure into
a graph, we constrain the encoder's latent space to be organised by event. We
regard this as the closest prior work and will compare against it directly.

\textbf{I-JEPA} \cite{assran2023ijepa} predicts one view's latent representation
from another rather than reconstructing raw content; this motivates our
label-efficient pretraining extension, where the ``views'' are two outlets'
reports of the same event.

\textbf{DiverseSumm} \cite{huang2024diversesumm} proposes QA-based evaluation of
whether a multi-document summary retains information from its sources.

\section{Contributions}

\begin{enumerate}\itemsep2pt
\item \textbf{An event-aware triplet objective} for multi-source neutral
summarization, in which the positive shares an event with the anchor and the
negative shares a \emph{publisher}. The shared publisher is the mechanism: it is
the only configuration in which satisfying the margin requires the encoder to
stop relying on outlet style. This addresses the first shortcoming --- that
standard training leaves the encoder's organising principle unconstrained ---
and a first run already shows a $63\%$ relative cut in publisher recoverability
at unchanged ROUGE.

\item \textbf{A length-normalised framing-bias metric}, together with the
evidence that the standard metric is a length proxy. We will report both, and
show what conclusions change. This addresses the third shortcoming, and makes
neutrality-by-deletion visible rather than rewarded.

\item \textbf{Evidence that representation leakage and output framing bias are
separable.} Our first experiment removes publisher information from the encoder
without moving the framing-bias metric at all, on identical output lengths.
Whether that gap can be closed --- and what it says about representation-level
interventions for surface-level bias --- is the question the rest of the project
answers.

\item \textbf{A publisher-disjoint evaluation protocol} for this corpus: a
stance-stratified split in which no outlet appears on two sides, plus a leakage
probe measured on both the standard and the disjoint split. This addresses the
fourth shortcoming --- that the released split cannot ask the generalisation
question.

\item \textbf{A reproducible data pipeline.} The released NeuS split files have
their publisher and stance labels stripped. We recover them at $100\%$ coverage
by normalising headlines against the union of all three crawl files --- the
obvious single-file join silently misses $13\%$ of examples. We will release
this alongside our code.
\end{enumerate}

\section{Proposed Methodology}

\paragraph{Backbone.} A pretrained sequence-to-sequence model (BART-large). The
$k$ articles for an event are concatenated into a single source sequence and
decoded into a neutral summary under the usual token-level cross-entropy loss
against the human-written roundup.

\paragraph{Event-aware triplet objective.} For an anchor article $a$, we sample a
positive $p$ --- a different outlet's report of the same event --- and a negative
$n$ --- the same outlet's report of a different event. Each is passed through the
shared encoder and pooled at the EOS position to give $f(\cdot)$. The training
signal is
\[
\mathcal{L} \;=\; \mathcal{L}_{\text{gen}} \;+\; \lambda \cdot
\max\!\big(0,\; \|f(a)-f(p)\| - \|f(a)-f(n)\| + m \big),
\]
with margin $m$ and weight $\lambda$ swept on the validation set. Because $a$ and
$n$ always share a publisher, the only way to satisfy the margin consistently is
to stop encoding publisher style and organise around the event instead.

Structurally this attaches at the same point as NeuS's own auxiliary-task hook,
which adds a weighted auxiliary loss over a pooled encoder state; we replace
their binary classification head with the triplet term and reimplement on a
current transformers stack rather than their pinned 2021 fork.

\paragraph{Sampling and training details.} Triples are mined offline: for each of
the 9{,}282 articles we draw two positives and two negatives, discarding anchors
whose outlet appears in only one story, which yields 18{,}368 triples. The
triplet term is fed as a second stream alongside the generation batch, mirroring
the reference implementation's separate auxiliary-input path, so the two losses
need not share a batch size. We start from the hyperparameters NeuS reports ---
learning rate $3\!\times\!10^{-5}$, effective batch 16, source length 512, target
length 250 --- and sweep only $\lambda$ and $m$. Training is short: 2{,}452
examples over five epochs is roughly 3{,}000 optimiser steps, under an hour of
single-GPU time per configuration, which is what makes the ablation grid and the
multi-seed reporting affordable.

\paragraph{Extension: cross-source pretraining.} Following the JEPA principle, we
pretrain by predicting one outlet's representation of an event from another
outlet's representation of the same event, using unlabelled same-event pairs from
WCEP-10 --- which needs neither publisher labels nor neutral targets, and so
suits this objective exactly. The question is whether the decoder can then be
fine-tuned on a fraction of the labelled summaries and still match a fully
supervised model. This extension is contingent on compute and will be dropped
before the core result if necessary.

\section{Dataset}

Our primary corpus is the NeuS release \cite{lee2022neus}: 3{,}095 AllSides
stories, each pairing three articles (one left, one centre, one right) with a
human-written neutral roundup, spanning 2016--2021 and 250 distinct outlets. It
is the only corpus available to us that carries an event grouping, a publisher
label and a neutral target on the same example. From it we derive 9{,}282
publisher-labelled articles and 18{,}368 training triples; 153 of the 250 outlets
appear in at least two distinct stories, which makes 9{,}188 of 9{,}285 article
slots usable as a same-publisher negative.

Two properties we flag honestly. First, AllSides publishes an \emph{excerpt} per
outlet rather than the full article --- a hard cap of 95 words, median 85 --- so
a three-source input is $\approx$241 words and fits comfortably within BART's
1{,}024-token window, but there is correspondingly little event content to
retain. Second, at 2{,}452 training examples the corpus is small for fine-tuning
BART-large; we will report mean and spread over at least three seeds rather than
single runs.

We use two further corpora in supporting roles. \textbf{WCEP-10}
\cite{ghalandari2020wcep} provides 10{,}200 genuine multi-document event clusters
of $\approx$9 articles each --- far better event structure and scale than NeuS
--- but its preprocessing stripped the publisher labels and its 25-word targets
were never written to be neutral, so we use it only for the pretraining
extension. \textbf{AllSides / Article-Bias-Prediction} \cite{baly2020detect}
supplies 37{,}554 publisher-labelled articles for additional negatives and for
its prebuilt media-based split. We note that its bias labels are outlet-level:
only 7 of its 491 sources ever carry more than one label, so ideology prediction
and publisher prediction are near-identical tasks there and we will not report
them as independent axes.

\section{Evaluation Plan}

\paragraph{Summary quality.} ROUGE-1/2/L against the human roundups, alongside
two training-free baselines the model must beat: concatenating all three source
articles, and copying the centrist source. The latter matters --- on the raw
framing-bias metric, copying the centre article already scores below both
partisan articles at zero training cost.

\paragraph{Framing bias, both ways.} The NeuS metric as published, and our
length-normalised variant (arousal per calibrated affect-bearing token), reported
side by side with mean generation length. A run that improves the raw score
without improving the normalised rate has learned to say less, and we will label
it as such. We additionally track MPQA strong-subjectivity density
\cite{wilson2005mpqa}, the one lexical signal in our audit that is \emph{not}
flat under length ($0.0633$/$0.0557$/$0.0631$ per calibrated token for
left/centre/right).

\paragraph{Publisher leakage.} A logistic probe on frozen encoder
representations, predicting the source outlet; lower is better. Measured on the
random split \emph{and} on our stance-stratified publisher-disjoint split, with a
TF-IDF probe as the lexical floor so that we can tell a representation-level
improvement from a vocabulary-level one. The comparison of interest is
generation-only training versus triplet training at matched ROUGE.

\paragraph{Event retention.} A DiverseSumm-style QA check
\cite{huang2024diversesumm}, scoped down honestly: with 85-word source excerpts
there is limited room for questions answerable from one source and not another,
so we treat this as a secondary content-preservation check rather than a headline
result.

\paragraph{Ablations.} $\lambda = 0$ (generation only); triplet with negatives
sampled from a \emph{different} publisher, which removes the adversarial pressure
and should recover most of the leakage; and a publisher-shuffled control.

\section{Feasibility and Risks}

The compute is modest: the full experimental grid --- baseline, triplet sweep,
ablations, and reruns on the disjoint split --- is roughly 20--25 single-GPU
hours at 16\,GB, which our cluster allocation covers comfortably. The data
pipeline, the metric, and the probes are CPU-only and already run.

Three risks we are managing explicitly. \textbf{Small data:} 2{,}452 training
examples will produce noticeable seed variance, so every reported number will be
a mean over at least three seeds with spread. \textbf{Weak event content:} the
85-word source excerpts limit what an event-retention evaluation can claim, so we
have demoted the QA check to a secondary measure rather than build the argument
on it. \textbf{A null result is possible:} the triplet term may reduce publisher
recoverability without improving summary neutrality, or may trade ROUGE for it.
We regard that as a reportable outcome --- combined with the metric critique it
still constitutes a complete finding about what representation-level constraints
can and cannot buy in this task.

\section{Timeline}

\begin{tabular}{@{}ll@{}}
\toprule
\textbf{Weeks} & \textbf{Milestone} \\
\midrule
1--2 & Data pipeline, splits, probe protocol \\
     & \emph{(complete)} \\
3--4 & NeuS baseline reproduced; metric \\
     & analysis written up \\
5--7 & Triplet objective; $\lambda$/margin sweep; \\
     & leakage probes on both splits \\
8--9 & Ablations, seed variance, QA check \\
10--11 & Cross-source pretraining extension \\
       & (compute permitting) \\
12 & Final report and code release \\
\bottomrule
\end{tabular}

\begin{thebibliography}{9}
\small

\bibitem[Assran et al.(2023)]{assran2023ijepa}
M.~Assran, Q.~Duval, I.~Misra, P.~Bojanowski, P.~Vincent, M.~Rabbat,
Y.~LeCun, and N.~Ballas.
\newblock Self-Supervised Learning from Images with a Joint-Embedding
Predictive Architecture.
\newblock In \emph{CVPR}, 2023.

\bibitem[Baly et al.(2020)]{baly2020detect}
R.~Baly, G.~Da San Martino, J.~Glass, and P.~Nakov.
\newblock We Can Detect Your Bias: Predicting the Political Ideology of News
Articles.
\newblock In \emph{EMNLP}, pages 4982--4991, 2020.

\bibitem[Ghalandari et al.(2020)]{ghalandari2020wcep}
D.~Gholipour Ghalandari, C.~Hokamp, N.~The Pham, J.~Glover, and G.~Ifrim.
\newblock A Large-Scale Multi-Document Summarization Dataset from the Wikipedia
Current Events Portal.
\newblock In \emph{ACL}, 2020.

\bibitem[Huang et al.(2024)]{huang2024diversesumm}
H.~Huang, Y.~Geng, et al.
\newblock Embrace Divergence for Richer Insights: A Multi-document
Summarization Benchmark and a Case Study on Summarizing Diverse Information
from News Articles.
\newblock In \emph{NAACL}, 2024.

\bibitem[Lee et al.(2022)]{lee2022neus}
N.~Lee, Y.~Bang, T.~Yu, A.~Madotto, and P.~Fung.
\newblock NeuS: Neutral Multi-News Summarization for Mitigating Framing Bias.
\newblock In \emph{NAACL}, 2022.

\bibitem[Lei and Huang(2024)]{lei2024multidoc}
Y.~Lei and R.~Huang.
\newblock Multi-document Summarization through Multi-document Event Relation
Graph Reasoning in LLMs: a Case Study in Framing Bias Mitigation.
\newblock In \emph{ACL}, 2024.

\bibitem[Lewis et al.(2020)]{lewis2020bart}
M.~Lewis, Y.~Liu, N.~Goyal, M.~Ghazvininejad, A.~Mohamed, O.~Levy,
V.~Stoyanov, and L.~Zettlemoyer.
\newblock BART: Denoising Sequence-to-Sequence Pre-training for Natural
Language Generation, Translation, and Comprehension.
\newblock In \emph{ACL}, 2020.

\bibitem[Mohammad(2018)]{mohammad2018vad}
S.~Mohammad.
\newblock Obtaining Reliable Human Ratings of Valence, Arousal, and Dominance
for 20,000 English Words.
\newblock In \emph{ACL}, 2018.

\bibitem[Wilson et al.(2005)]{wilson2005mpqa}
T.~Wilson, J.~Wiebe, and P.~Hoffmann.
\newblock Recognizing Contextual Polarity in Phrase-Level Sentiment Analysis.
\newblock In \emph{HLT/EMNLP}, 2005.

\end{thebibliography}

\end{document}
